% Engineering design (BGU \S 10.7.2.6).
\section{Engineering Design}
\label{sec:design}

This section documents the chosen approach in detail. The system block
diagram is shown in Figure \ref{fig:system}; the degradation pipeline that
drives the experimental program is shown in Figure \ref{fig:pipeline}.

\begin{figure}[htbp]
  \centering
  \includegraphics[width=0.95\textwidth]{\figpath{block_diagrams/system_overview.png}}
  \caption{System block diagram. An input image from CIFAR-10 or MNIST passes through the THz-like degradation pipeline, is normalized to the ImageNet distribution, is fed to one of the three backbones, and is classified by a global-average-pooling head.}
  \label{fig:system}
\end{figure}

\begin{figure}[htbp]
  \centering
  \includegraphics[width=0.95\textwidth]{\figpath{block_diagrams/pipeline_overview.png}}
  \caption{Pipeline version 2. The saturation lerp runs first, so grayscale luminance is computed from the un-degraded source; noise and salt-and-pepper act at the pre-upsample resolution, so the bicubic upsample spreads each sample into sensor-grain structure.}
  \label{fig:pipeline}
\end{figure}

\subsection{Mathematical model}
\label{sec:design:math}

\textbf{Degradation operator.}\quad
Let $x \in [0, 1]^{3 \times H \times W}$ be a clean image and let
$\phi = (s, r, \sigma_n, K_b, \sigma_b, p)$ denote the degradation
parameters: saturation $s \in [0, 1]$, low-resolution target $r \in \mathbb{N}$,
noise standard deviation $\sigma_n \geq 0$, blur kernel size $K_b$, blur
sigma $\sigma_b$, and salt-and-pepper amount $p \in [0, 1]$. Define the
composite operator
\begin{equation}
  \mathcal{D}_{\phi}(x) = \mathrm{Clip}_{[0,1]} \circ \mathrm{Blur}_{K_b, \sigma_b} \circ \mathrm{Up}_{224} \circ \mathrm{SP}_p \circ \mathrm{N}_{\sigma_n} \circ \mathrm{Down}_r \circ \mathrm{Sat}_s,
  \label{eq:D}
\end{equation}
where $\mathrm{Sat}_s(x) = (1 - s)\, \mathrm{Gray}(x) + s \, x$ with
$\mathrm{Gray}(x) = 0.2989\,x_R + 0.5870\,x_G + 0.1140\,x_B$ (BT.601 luma);
$\mathrm{Down}_r$ is bilinear downsample to $r \times r$;
$\mathrm{N}_{\sigma_n}$ adds independent Gaussian noise of standard
deviation $\sigma_n$; $\mathrm{SP}_p$ flips a $p/2$ fraction of pixels to 1
and another $p/2$ fraction to 0; $\mathrm{Up}_{224}$ is bicubic upsample
to $224 \times 224$; and $\mathrm{Blur}_{K_b, \sigma_b}$ is a separable
Gaussian blur with kernel size $K_b$ and standard deviation $\sigma_b$.

\textbf{Classifier.}\quad
A backbone $f_{\theta}(\cdot)$ produces a logit vector
$z = f_{\theta}(\mathcal{D}_{\phi}(x))$. With label smoothing $\alpha = 0.1$,
the training objective is the smoothed cross-entropy
\begin{equation}
  \mathcal{L}(\theta) = -\sum_{c=1}^{C} \tilde y_c \log \mathrm{softmax}(z)_c,
  \qquad
  \tilde y_c = (1 - \alpha)\, \mathbf{1}[c = y] + \frac{\alpha}{C}.
  \label{eq:loss}
\end{equation}

% V4: removed the three off-the-shelf architectural equations (ResNet
% identity skip, DenseNet concatenation, TransNeXt foveal attention).
% They are standard, fully documented in the cited primary papers, and
% not part of the methodological novelty of this work.
\textbf{Backbones.}\quad The three architectures expose different
inductive biases for the low-resolution regime. ResNet50
[1] combines features via additive identity skips, which
reuse previous-layer features but collapse them into the running
representation; subtle low-level cues can be washed out when the later
residual blocks dominate. DenseNet121 [2]
concatenates the feature maps of every preceding layer within each
dense block, explicitly preserving the collective knowledge of all
earlier layers --- the property our preliminary hypothesis expected to
help under information loss. TransNeXt-tiny [3]
uses pixel-focused attention with parallel foveal and peripheral
streams: a window-shaped mask down-weights all but the nearest
neighbours of each query, while a global pooling stream contributes a
coarse-grained peripheral signal. The biomimetic prior is designed to
retain accuracy when the global shape is corrupted but local texture
is partially preserved. The architectural equations are standard and
are stated in full in the cited primary papers; we omit them here.

\subsection{Image-quality metrics}
\label{sec:design:psnr}

Each degradation pipeline is anchored to a $\mathrm{PSNR}$ and $\mathrm{SSIM}$
probe so that cross-phase comparisons sit on a physical scale.

\textbf{PSNR} [4] between a reference image $x$ and a
degraded image $\tilde x = \mathcal{D}_{\phi}(x)$, both in $[0, 1]^{3 \times 224 \times 224}$, is
\begin{equation}
  \mathrm{PSNR}(x, \tilde x) = 10 \log_{10} \frac{1}{\mathrm{MSE}(x, \tilde x)},
  \qquad
  \mathrm{MSE}(x, \tilde x) = \frac{1}{3 \cdot 224^2} \sum_{c, i, j} (x_{c,i,j} - \tilde x_{c,i,j})^2.
  \label{eq:psnr}
\end{equation}

\textbf{SSIM} [5] is the structural similarity index. For
two local patches $u, v$ with means $\mu_u, \mu_v$, variances
$\sigma_u^2, \sigma_v^2$, and covariance $\sigma_{uv}$,
\begin{equation}
  \mathrm{SSIM}(u, v) = \frac{(2 \mu_u \mu_v + C_1)(2 \sigma_{uv} + C_2)}{(\mu_u^2 + \mu_v^2 + C_1)(\sigma_u^2 + \sigma_v^2 + C_2)},
  \label{eq:ssim}
\end{equation}
with $C_1, C_2$ small stabilizers; the full-image score is the spatial mean
over an $11 \times 11$ Gaussian window. The implementation uses
\texttt{torchmetrics.image.StructuralSimilarityIndexMeasure} with
\texttt{data\_range = 1}.

\textbf{Reference construction.}\quad The PSNR / SSIM reference is the
input image upsampled to $224 \times 224$ by bilinear interpolation. The
choice matches what a hypothetical clean-image classifier would see if the
degradation pipeline were short-circuited; it does not include the
saturation lerp, so the measurements isolate the spatial loss from the
color loss. Both choices are documented in the figure captions in
Section \ref{sec:tests} so the reader can interpret the absolute scale
correctly.

\textbf{Probe protocol.}\quad For each pipeline the probe iterates over
the first 256 validation indices, applies the pipeline with seed offset
$2 \times 10^7$ (chosen disjoint from the training and validation streams),
and aggregates the mean and standard deviation. The 256-sample size
balances probe runtime against the standard error of the mean (below
$0.1$ dB on the L3 pipelines).

\subsection{Optimizer and learning-rate schedule}
\label{sec:design:optim}

All cells train with AdamW [9] under decoupled
weight decay. With $g_t = \nabla_\theta \mathcal{L}(\theta_{t-1})$ the
loss gradient at step $t$, the running first and second moment estimates
are
\begin{equation}
  m_t = \beta_1 m_{t-1} + (1 - \beta_1) g_t,
  \qquad
  v_t = \beta_2 v_{t-1} + (1 - \beta_2) g_t^{2},
  \label{eq:adamw_moments}
\end{equation}
and the parameter update is
\begin{equation}
  \theta_t = \theta_{t-1} - \eta_t \!\left(
    \frac{m_t / (1 - \beta_1^{t})}{\sqrt{v_t / (1 - \beta_2^{t})} + \epsilon}
    + \lambda \, \theta_{t-1}
  \right),
  \label{eq:adamw}
\end{equation}
with $\beta_1 = 0.9$, $\beta_2 = 0.999$, $\epsilon = 10^{-8}$, and weight
decay $\lambda \in \{10^{-4}, 5 \times 10^{-2}\}$ (CNN / TransNeXt).
Gradient clipping caps the global norm at $1.0$.

The learning rate $\eta_t$ follows a cosine annealing schedule
[10] between $\eta_{\max}$ and $\eta_{\min}$ over a
total budget of $T$ training steps:
\begin{equation}
  \eta_t = \eta_{\min} + \tfrac{1}{2}\, (\eta_{\max} - \eta_{\min}) \!\left(
    1 + \cos\!\left( \frac{t}{T}\, \pi \right)
  \right).
  \label{eq:cosine}
\end{equation}
Differential learning rates apply $\eta_{\max}$ separately to the
classifier head ($10^{-3}$ on CNNs, $5 \times 10^{-4}$ on TransNeXt-tiny)
and to the backbone parameters ($10^{-4}$ on CNNs, $5 \times 10^{-5}$ on
TransNeXt-tiny). Training uses bf16-mixed precision on the RTX 5070 host
and runs for at most 60 epochs with early stopping at patience 10 on
validation accuracy.

\subsection{Calibration metric}
\label{sec:design:ece}

Expected calibration error (ECE) [11], [12]
bins the predicted probabilities into $M = 15$ equal-width bins and computes
\begin{equation}
  \mathrm{ECE} = \sum_{m=1}^{M} \frac{|B_m|}{N} \left| \mathrm{acc}(B_m) - \mathrm{conf}(B_m) \right|,
  \label{eq:ece}
\end{equation}
where $B_m$ is the set of predictions whose top-1 probability falls in bin
$m$, $\mathrm{acc}(B_m)$ is the empirical accuracy of $B_m$, and
$\mathrm{conf}(B_m)$ is the mean predicted probability in $B_m$. A
well-calibrated model has $\mathrm{ECE} \to 0$.

\subsection{Regularization treatments (Phase D)}
\label{sec:design:phase_d}

Phase D layers three regularization treatments on top of the frozen
Phase B Optuna winners:

\textbf{T1.}\quad Dropout (CNN, rate 0.2) or stochastic-depth drop-path
[13] (TransNeXt, rate 0.2). Architectural regularization
only; no label-side intervention.

\textbf{T2.}\quad mixup [14] with $\alpha = 0.2$;
label-mixing only.

\textbf{T3.}\quad The combination: dropout (or drop-path)
$+$ mixup $\alpha = 0.2$ $+$ CutMix [15] $\alpha = 1.0$.
The kitchen-sink combination defines the upper bound of training-side
recovery measurable inside the design space.

\subsection{Software architecture}
\label{sec:design:sw}

The implementation is layered:

\begin{itemize}
  \item \texttt{src/data/} hosts the single source of truth for the
        degradation curve (\texttt{degradation\_levels.py}) and the operator
        (\texttt{degrade.py}). Both modules are imported verbatim into the
        training and the image-quality probe paths, which is the property
        that guarantees the measured PSNR / SSIM correspond to the pixels
        the model actually sees.
  \item \texttt{src/lightning/} hosts the training engine
        (\texttt{THzClassifier}, \texttt{THzDataModule}) built on PyTorch
        Lightning [16].
  \item \texttt{src/models/} hosts the backbone wrappers; the TransNeXt
        weights are auto-downloaded from the public checkpoint and cached
        under \texttt{artifacts/weights/}.
  \item \texttt{scripts/} hosts the campaign drivers
        (\texttt{run\_all\_phases.py}) and the diagnostics builders.
  \item \texttt{Final\_Report/V6/scripts/} hosts the figure / table / sample
        generators and the DOCX builder that produce this report. Every
        artifact under \texttt{Final\_Report/figures/v6/} and
        \texttt{Final\_Report/V6/} is reproducible by running
        \texttt{python Final\_Report/V6/scripts/build\_v6\_*.py} from a
        clean checkout.
\end{itemize}
